% CC-ResDiff: Global Color-Consistency Supervision for Residual-Space Diffusion
% Super-Resolution.
%
% Written against the AAAI-style presentation of ResDiff (Shang et al. 2024);
% swap \documentclass and the bibliography for whichever template you submit to.
%
% NOTE: numbers marked [PENDING] were still training when this draft was
% generated. Do not submit without replacing them.

\documentclass[11pt]{article}
\usepackage[margin=1in]{geometry}
\usepackage{amsmath,amssymb}
\usepackage{booktabs}
\usepackage{graphicx}
\usepackage{multirow}
\usepackage{xcolor}
\usepackage[hidelinks]{hyperref}

\newcommand{\pending}[1]{\textcolor{red}{[PENDING: #1]}}
\newcommand{\best}[1]{\textbf{#1}}

\title{\textbf{CC-ResDiff: Global Color-Consistency Supervision for\\
Residual-Space Diffusion Super-Resolution}}
\author{Sudhir Kumar}
\date{}

\begin{document}
\maketitle

\begin{abstract}
Residual-space diffusion models for single image super-resolution (SISR), such
as SRDiff and ResDiff, train a diffusion model to predict the residual between a
CNN's initial prediction and the ground truth. Because the residual is dominated
by high-frequency content, and because recent work has focused on strengthening
high-frequency guidance still further, the diffusion branch receives almost no
supervision on the \emph{global} low-frequency attributes of its output. ResDiff
explicitly identifies the consequence---a large color discrepancy when the model
is under-trained---and proposes ``utilizing a global color feature'' as future
work. We implement and study that proposal. We introduce a Global
Color-Consistency (GCC) loss: an auxiliary term that penalises the difference
between spatially pooled color statistics of the model's implied clean prediction
and those of the ground truth. The loss is training-time only, adds no
parameters, and leaves inference cost identical to the baseline. On a
compute-constrained CelebA benchmark the GCC loss reduces mean color error by
$46$--$71\%$ across two independent training runs and improves LR-consistency by
$0.69$\,dB, at no measurable distortion cost. We further show through a
gradient-level analysis that the term's influence is not uniform over the
diffusion trajectory but concentrated at large timesteps, spanning five orders of
magnitude, and that this skew is what drives an accompanying FID regression. We
report a negative result on DIV2K, where the baseline exhibits little color drift
and the loss yields no benefit, and conclude that GCC supervision is best
understood as a targeted remedy for a specific and identifiable failure mode
rather than a universal improvement.
\end{abstract}

\section{Introduction}

Single image super-resolution is ill-posed: many high-resolution (HR) images are
consistent with a given low-resolution (LR) input. Diffusion probabilistic models
have become a leading approach, and residual-space variants have proven
particularly effective. SRDiff~\cite{li2022srdiff} and
ResDiff~\cite{shang2024resdiff} first use a CNN to recover the primary
low-frequency content, then train a diffusion model to generate only the
\emph{residual} between that prediction and the ground truth. Restricting the
diffusion model to the residual accelerates convergence and improves sample
quality relative to modelling the image directly.

This design has an under-examined consequence. The residual is, by construction,
mostly high-frequency; the standard $\epsilon$-prediction objective supervises
the noise added to it and therefore says little about the absolute color of the
reconstruction. ResDiff sharpens this further by adding frequency-domain losses
and high-frequency-guided cross-attention---improving detail while supplying no
additional color signal. The authors observe the resulting failure directly,
noting that ``the color will appear a large discrepancy when the model is
under-trained,'' attributing it to ``a lack of color features in the guided
high-frequency information,'' and proposing that ``utilizing a global color
feature may well address this issue in future work''~\cite{shang2024resdiff}.

We take up that proposal. Our contributions are:

\begin{itemize}
\item \textbf{A Global Color-Consistency (GCC) loss} for residual-space diffusion
SISR: an auxiliary training-time term that supervises spatially pooled color
statistics of the model's implied clean prediction. It adds no parameters and
does not alter the sampling procedure, so inference cost is unchanged.

\item \textbf{A color-error metric} for SISR that measures precisely the
phenomenon the standard metric suite misses. We show empirically that PSNR, SSIM
and LPIPS can improve while color error worsens.

\item \textbf{A gradient-level analysis} showing that the auxiliary term's
influence is highly non-uniform across the diffusion trajectory---$0.002\%$ of
the primary gradient at $t{=}0$ versus $236\%$ at $t{=}T{-}1$. Loss magnitude is
a misleading proxy here, understating the term's true influence by roughly
$500\times$. This analysis explains both the observed FID regression and the
existence of an optimal loss weight.

\item \textbf{An honest characterisation of scope}, including a negative result:
on DIV2K, where the baseline's color error is already five times lower, the loss
provides no benefit. We argue the method should be understood as conditional on
the presence of drift.
\end{itemize}

\section{Related Work}

\paragraph{Diffusion-based SISR.}
SR3~\cite{saharia2022image} adapts DDPM to SISR by conditioning on a bicubically
upsampled LR image. SRDiff~\cite{li2022srdiff} introduces residual prediction,
training the diffusion model on the difference between the ground truth and an
RRDB prediction. ResDiff~\cite{shang2024resdiff} extends this with a
frequency-domain information splitter, high-frequency-guided cross-attention, and
FFT/DWT losses on the CNN, reporting faster convergence and better sample
quality.

\paragraph{Color consistency in image restoration.}
Color shift is a known artefact in generative restoration. Classical remedies are
post-hoc: histogram matching and Reinhard-style color transfer correct output
statistics after generation. Such corrections are cheap but operate outside the
model and cannot influence what it learns. We compare against this alternative
directly (Section~\ref{sec:colortransfer}).

\paragraph{Loss weighting in diffusion models.}
Reweighting the diffusion objective across timesteps is well studied for the
primary loss. Our analysis shows that an \emph{auxiliary} loss defined on the
reconstructed $x_0$ inherits an extreme implicit timestep weighting from the
$1/\sqrt{\bar\alpha_t}$ factor in that reconstruction, an effect that appears not
to have been noted previously and which we exploit as a control.

\section{Method}

\subsection{Background}

Let $x_{LR}$ be the low-resolution input, $x_{HR}$ the ground truth, and
$x_{up}$ the bicubically upsampled LR image. A pre-trained CNN produces
$x_{cnn}$. Residual-space diffusion defines the target as
\begin{equation}
r = s\,(x_{HR} - x_{up}),
\end{equation}
with a rescaling constant $s$ chosen so $r$ occupies $[-1,1]$. The forward
process is standard:
\begin{equation}
r_t = \sqrt{\bar\alpha_t}\,r + \sqrt{1-\bar\alpha_t}\,\epsilon,
\quad \epsilon \sim \mathcal{N}(0, I),
\end{equation}
and the model $\epsilon_\theta(r_t, t, c)$ is trained with
\begin{equation}
\mathcal{L}_{\text{DDPM}} = \mathbb{E}\big[\,\|\epsilon - \epsilon_\theta(r_t,t,c)\|_1\,\big],
\label{eq:ddpm}
\end{equation}
where $c$ denotes CNN-derived conditioning features. Nothing in
Eq.~\ref{eq:ddpm} references the absolute color of the reconstruction.

\subsection{Global Color-Consistency Loss}

At each training step we recover the model's implied clean residual using the
standard closed form,
\begin{equation}
\hat r_0 = \frac{r_t - \sqrt{1-\bar\alpha_t}\;\epsilon_\theta(r_t,t,c)}{\sqrt{\bar\alpha_t}},
\label{eq:x0}
\end{equation}
clamp it to $[-1,1]$ (the range in which the target residual is constructed), and
map it back to image space, $\hat x_{HR} = \hat r_0 / s + x_{up}$. We then compare
spatially pooled color statistics:
\begin{equation}
\mathcal{L}_{\text{GCC}} =
\big\| \, \mathcal{P}_k(\hat x_{HR}) - \mathcal{P}_k(x_{HR}) \, \big\|_2^2 ,
\label{eq:gcc}
\end{equation}
where $\mathcal{P}_k$ is adaptive average pooling to a $k \times k$ grid. Pooling
is what makes the term complementary rather than redundant: at $k=8$ it discards
texture and edges, retaining only the spatial color balance that
Eq.~\ref{eq:ddpm} leaves unconstrained. The total objective is
\begin{equation}
\mathcal{L} = \mathcal{L}_{\text{DDPM}} + \lambda\,\mathcal{L}_{\text{GCC}}.
\end{equation}

The clamp in Eq.~\ref{eq:x0} is essential, not cosmetic. Under a cosine schedule
$\bar\alpha_t$ reaches $2.4\times10^{-7}$ at $t = T-1$, so the reconstruction
amplifies by ${\sim}2000\times$; unclamped, we measured
$\mathcal{L}_{\text{GCC}} = 2.8\times10^{4}$ at $t=T-1$ against a DDPM loss of
${\sim}0.8$. Since $t$ is sampled uniformly, roughly one step in a hundred would
destabilise training.

\subsection{Timestep Weighting}
\label{sec:weighting}

The clamp bounds the loss value but not its gradient. Because
Eq.~\ref{eq:x0} carries the $1/\sqrt{\bar\alpha_t}$ factor, the term's effective
gradient scale at $t{=}T{-}1$ is ${\sim}2029\times$ that at $t{=}0$. We therefore
consider a reweighted variant
\begin{equation}
\mathcal{L}_{\text{GCC}}^{w} = \frac{\sum_i w_{t_i}\,\ell_i}{\sum_i w_{t_i}},
\qquad w_t = \bar\alpha_t ,
\end{equation}
which cancels the amplification and concentrates supervision at small $t$, where
the clean-image estimate is meaningful. Section~\ref{sec:snr} evaluates whether
this removes the FID regression.

\section{Experiments}

\subsection{Setup}

We evaluate in a deliberately compute-constrained regime, matching the setting in
which ResDiff reports color discrepancy to be most visible. CelebA is subsampled
to 3{,}000 training / 300 validation / 300 test images at $16\times16 \to
64\times64$ ($4\times$). Stage 1 is an RRDB with hidden size 16 and 4 blocks;
Stage 2 is a UNet with hidden size 32 and $2.28$M parameters, $T=100$ timesteps,
cosine schedule. Both stages train for up to 20{,}000 steps with early stopping
on PSNR. \textbf{Crucially, the baseline and proposed arms are identical in every
respect---data, schedule, seed, initialisation and Stage-1 checkpoint---except for
the presence of the GCC term.}

Early stopping monitors PSNR rather than color error: stopping on the metric the
method targets would bias model selection toward the proposed arm. Evaluation
loads the best checkpoint, not the final one.

\paragraph{Metrics.} PSNR, SSIM, LPIPS and FID, plus \textbf{color error}: the
mean absolute difference between per-channel RGB means of the generated and
ground-truth images, averaged over the test set. We also report LR-PSNR
(consistency between the downsampled output and the LR input), which measures
low-frequency fidelity and is therefore the quantity the GCC loss most directly
targets.

\paragraph{Reproducibility.} Evaluation is deterministic given a seed and both
arms are evaluated under identical seeds, so the comparison is paired on sampling
noise. We measured an evaluation noise floor by scoring a single fixed checkpoint
under four seeds (Table~\ref{tab:noise}).

\begin{table}[t]
\centering
\small
\begin{tabular}{lrrr}
\toprule
Metric & Mean & Std & Range \\
\midrule
PSNR        & 25.0850 & 0.0175 & 0.0386 \\
SSIM        & 0.8146  & 0.0004 & 0.0010 \\
LPIPS       & 0.0461  & 0.0005 & 0.0011 \\
LR-PSNR     & 43.7005 & 0.0093 & 0.0221 \\
Color error & 0.8229  & 0.0022 & 0.0050 \\
FID         & 69.3892 & 0.1525 & 0.3564 \\
\bottomrule
\end{tabular}
\caption{Evaluation noise floor: one fixed checkpoint scored under four seeds.
This bounds \emph{evaluation} variance only; training-seed variance is larger and
is reported separately in Table~\ref{tab:main}.}
\label{tab:noise}
\end{table}

\subsection{Main Results}

Table~\ref{tab:main} reports two independent training runs on different hardware.

\begin{table}[t]
\centering
\small
\begin{tabular}{lccc}
\toprule
Metric & Baseline & CC-ResDiff & $\Delta$ \\
\midrule
Color error $\downarrow$ & $0.606 \pm 0.308$ & \best{$0.224 \pm 0.026$} & $-63\%$ \\
LR-PSNR $\uparrow$       & $44.42 \pm 1.03$  & \best{$45.11 \pm 0.85$}  & $+0.69$ \\
LPIPS $\downarrow$       & $0.0452 \pm 0.0020$ & \best{$0.0440 \pm 0.0027$} & $-0.0012$ \\
PSNR $\uparrow$          & \best{$25.094 \pm 0.001$} & $25.040 \pm 0.103$ & $-0.055$ \\
SSIM $\uparrow$          & \best{$0.8155 \pm 0.0016$} & $0.8144 \pm 0.0029$ & $-0.0012$ \\
FID $\downarrow$         & \best{$67.43 \pm 2.67$} & $69.65 \pm 5.05$ & $+2.22$ \\
\bottomrule
\end{tabular}
\caption{CelebA $4\times$, mean $\pm$ std over two independent training runs.
The color-error reduction is the largest and most consistent effect
($-70.6\%$ and $-46.9\%$ individually). The PSNR difference changes sign between
runs and is therefore not established.}
\label{tab:main}
\end{table}

Three observations. First, the two quantities the loss directly
supervises---color error and LR-PSNR---improve consistently and by margins far
exceeding the evaluation noise floor. Second, the apparent PSNR cost is not
reproducible: it is $-0.13$\,dB in one run and $+0.02$\,dB in the other, so at
this sample size we cannot claim a distortion penalty. Third, FID is worse in
both runs, though its magnitude varies sevenfold ($+3.90$, $+0.54$); we treat the
direction as real and the magnitude as unresolved.

\subsection{Comparison with Color Transfer}
\label{sec:colortransfer}

A natural objection is that post-hoc color correction might achieve the same
effect for free. We apply per-channel mean and mean/standard-deviation
(Reinhard-style) matching to the baseline's outputs, using the bicubically
upsampled LR image as reference---the only color reference available at inference
time. Matching to the ground truth would trivially minimise our metric while
being unusable in practice, and is therefore inadmissible.

Neither variant helps (Table~\ref{tab:ct}). Both make color error \emph{worse},
and matching the standard deviation additionally destroys LR-consistency,
costing $7.4$\,dB of LR-PSNR: rescaling per-channel variance alters the image's
low-frequency content, which is exactly what LR-PSNR measures.

\begin{table}[t]
\centering
\small
\begin{tabular}{lccc}
\toprule
Method & ColorErr $\downarrow$ & LR-PSNR $\uparrow$ & PSNR $\uparrow$ \\
\midrule
Baseline                & 0.8234 & 43.690 & 25.095 \\
\; + mean transfer      & 0.8746 & 43.558 & 25.064 \\
\; + mean/std transfer  & 0.9333 & 36.243 & 24.915 \\
\midrule
\emph{reference floor} ($x_{up}$ vs $x_{HR}$) & \emph{0.4358} & -- & -- \\
\midrule
CC-ResDiff (ours)       & \best{0.2423} & \best{44.509} & 24.967 \\
\bottomrule
\end{tabular}
\caption{Post-hoc color transfer against learned supervision, CelebA, 300 test
images. Post-hoc correction is bounded below by the color error of the reference
it matches to; CC-ResDiff operates $1.8\times$ below that bound.}
\label{tab:ct}
\end{table}

The reason is structural rather than incidental. Any post-hoc correction can be
no more accurate than the reference it matches to, and the only reference
available at inference---the upsampled LR image---itself differs from the ground
truth by a color error of $0.4358$, measured over the same test set. That value
is a hard floor on what post-processing can achieve. CC-ResDiff reaches $0.2423$,
$1.8\times$ below the floor, because supervision is applied during training
against the true HR image; the learned model does not need a test-time reference
at all. This is the central argument for a learned solution over a corrective
one.

\subsection{Timestep-Reweighted Variant}
\label{sec:snr}

\pending{Results for $w_t = \bar\alpha_t$; run in progress. Tests whether
removing the high-$t$ skew eliminates the FID regression.}

\subsection{Ablations}

Table~\ref{tab:abl} sweeps $\lambda$ and the pooling size $k$ at a reduced
5{,}000-step budget, with an identical budget for every run so that the sweep
does not partly measure training length.

\begin{table}[t]
\centering
\small
\begin{tabular}{lrrrrr}
\toprule
Setting & ColorErr $\downarrow$ & PSNR $\uparrow$ & LPIPS $\downarrow$ & FID $\downarrow$ & LR-PSNR $\uparrow$ \\
\midrule
\multicolumn{6}{l}{\emph{$\lambda$ sweep ($k=8$)}}\\
none        & 0.4833 & 23.857 & \best{0.0612} & 111.51 & 41.445 \\
0.01        & 0.3382 & \best{23.955} & 0.0635 & \best{109.11} & \best{41.918} \\
0.1         & \best{0.1952} & 23.941 & 0.0653 & 110.43 & 41.865 \\
0.5         & 0.4985 & 23.721 & 0.0665 & 115.43 & 41.192 \\
1.0         & 0.2162 & 23.727 & 0.0680 & 116.71 & 40.976 \\
\midrule
\multicolumn{6}{l}{\emph{pooling size sweep ($\lambda=0.1$)}}\\
$k=4$       & 0.3385 & 23.839 & 0.0641 & 112.47 & 41.613 \\
$k=8$       & \best{0.1952} & \best{23.941} & 0.0653 & \best{110.43} & \best{41.865} \\
$k=16$      & 0.3522 & 23.919 & 0.0658 & 111.05 & 41.821 \\
\bottomrule
\end{tabular}
\caption{Ablations at a 5{,}000-step budget. Comparable to each other, not to
Table~\ref{tab:main}. Five of six GCC settings beat the no-GCC reference on color
error, and $k=8$ is best on every metric. LPIPS degrades monotonically in
$\lambda$.}
\label{tab:abl}
\end{table}

We deliberately do \emph{not} claim an optimal $\lambda$. Color error across the
six GCC runs has standard deviation $0.109$, comparable to the differences
between settings, and the $\lambda{=}0.5$ point is inconsistent with both
neighbours. The sweep supports a usable \emph{range} ($\lambda \in [0.01, 0.1]$,
$k=8$); resolving a dose-response curve would require repeated runs per setting.

\subsection{Gradient Analysis}

Loss values suggest the GCC term is negligible---$0.06\%$ of the DDPM loss at
$\lambda=0.1$. Measuring parameter gradients on a trained checkpoint tells a
different story (Table~\ref{tab:grad}): the ratio is $33.9\%$ averaged over
uniformly sampled $t$, roughly $500\times$ larger than the loss ratio implies.

\begin{table}[t]
\centering
\small
\begin{tabular}{lrrr}
\toprule
$t$ & $\mathcal{L}_{\text{GCC}}$ & clamped & $\lambda\|\nabla \mathcal{L}_{\text{GCC}}\| / \|\nabla \mathcal{L}_{\text{DDPM}}\|$ \\
\midrule
0        & $1.2\times10^{-5}$ & 0.4\%  & $0.002\%$ \\
25       & $2.8\times10^{-5}$ & 0.2\%  & $0.080\%$ \\
50       & $5.4\times10^{-5}$ & 0.1\%  & $0.398\%$ \\
75       & $8.4\times10^{-5}$ & 0.1\%  & $1.276\%$ \\
99       & $3.4\times10^{-2}$ & 97.4\% & $\mathbf{235.7\%}$ \\
uniform  & $9.8\times10^{-4}$ & 1.5\%  & $\mathbf{33.9\%}$ \\
\bottomrule
\end{tabular}
\caption{The GCC term's gradient contribution spans five orders of magnitude
across the trajectory. At $t{=}T{-}1$ nearly all elements saturate the clamp, yet
the remainder still carry the ${\sim}2000\times$ amplification.}
\label{tab:grad}
\end{table}

This has two consequences. It explains the ablation's cost structure: at
$\lambda=0.5$ the auxiliary gradient would reach ${\sim}170\%$ of the primary
one, so the denoising objective is crowded out and all metrics degrade. And it
predicts the FID regression: a term acting predominantly at large $t$ pulls global
tone toward the dataset mean, which suppresses the distributional diversity FID
rewards while leaving per-image LPIPS largely indifferent.

\subsection{Negative Result: DIV2K}

We repeated the experiment on DIV2K at matched scale
(Table~\ref{tab:div2k}). The result does not replicate.

\begin{table}[t]
\centering
\small
\begin{tabular}{lccc}
\toprule
Metric & Baseline & CC-ResDiff & $\Delta$ \\
\midrule
Color error $\downarrow$ & \best{0.1650} & 0.2036 & $+23.4\%$ \\
LR-PSNR $\uparrow$       & \best{47.248} & 47.101 & $-0.15$ \\
PSNR $\uparrow$          & 29.650 & \best{29.703} & $+0.05$ \\
SSIM $\uparrow$          & 0.7635 & \best{0.7646} & $+0.001$ \\
LPIPS $\downarrow$       & \best{0.2570} & 0.2630 & $+0.006$ \\
FID $\downarrow$         & \best{82.71} & 83.25 & $+0.54$ \\
\bottomrule
\end{tabular}
\caption{DIV2K $4\times$, 512 test tiles. Every sign flips relative to
Table~\ref{tab:main} except FID.}
\label{tab:div2k}
\end{table}

The most plausible explanation is headroom. The DIV2K baseline's color error is
$0.165$ against CelebA's $0.606$---roughly five times lower---so there is little
drift for the loss to correct, and it contributes a small regression instead. This
is consistent with ResDiff's own observation that the discrepancy appears when the
model is under-trained: the DIV2K Stage-1 network is comparatively strong here
(LR-PSNR $47.2$ versus $44.4$), leaving less for the diffusion branch to get
wrong.

We report this rather than omitting it because it defines the method's scope. GCC
supervision is a remedy for a specific failure mode, and its benefit is
conditional on that failure being present.

\section{Limitations}

\begin{itemize}
\item \textbf{Scale.} All experiments use 3{,}000 training images at
$64\times64$ with a $2.28$M-parameter diffusion model. Whether the effect persists
at the resolutions and model sizes used by SRDiff and ResDiff is untested.
\item \textbf{Sample size.} Two training runs per arm on CelebA, one on DIV2K.
The color-error effect is large enough to be robust; the PSNR and FID differences
are not resolved at this $n$.
\item \textbf{No DIV2K noise floor}, so the $+23.4\%$ regression there is not
established as statistically real.
\item \textbf{FID sample counts are small} (300 CelebA, 512 DIV2K), biasing FID
upward in absolute terms; only the relative comparison is meaningful.
\item \textbf{Drift is not predicted in advance.} We show the method helps where
color error is high, but offer no way to determine beforehand whether a given
training configuration will drift.
\end{itemize}

\section{Conclusion}

We implemented the global color feature that ResDiff identified as future work,
and characterised when it helps. The GCC loss is simple, parameter-free at
inference, and substantially reduces color drift in the regime where that drift
occurs. Our gradient analysis shows that an auxiliary loss defined on the
reconstructed $x_0$ inherits an extreme implicit timestep weighting---a
methodological point that applies to any such term, not only ours. The negative
DIV2K result bounds the claim: this is a targeted remedy, not a universal
improvement. Predicting which configurations will drift, and evaluating the
timestep-reweighted variant at scale, are the natural next steps.

\begin{thebibliography}{9}
\bibitem{shang2024resdiff}
S.~Shang, Z.~Shan, G.~Liu, L.~Wang, X.~Wang, Z.~Zhang, and J.~Zhang.
\newblock ResDiff: Combining CNN and Diffusion Model for Image Super-Resolution.
\newblock In {\em AAAI}, 2024.

\bibitem{li2022srdiff}
H.~Li, Y.~Yang, M.~Chang, S.~Chen, H.~Feng, Z.~Xu, Q.~Li, and Y.~Chen.
\newblock SRDiff: Single Image Super-Resolution with Diffusion Probabilistic Models.
\newblock {\em Neurocomputing}, 479:47--59, 2022.

\bibitem{saharia2022image}
C.~Saharia, J.~Ho, W.~Chan, T.~Salimans, D.~J. Fleet, and M.~Norouzi.
\newblock Image Super-Resolution via Iterative Refinement.
\newblock {\em IEEE TPAMI}, 2022.

\bibitem{ho2020denoising}
J.~Ho, A.~Jain, and P.~Abbeel.
\newblock Denoising Diffusion Probabilistic Models.
\newblock In {\em NeurIPS}, 2020.

\bibitem{wang2018esrgan}
X.~Wang, K.~Yu, S.~Wu, J.~Gu, Y.~Liu, C.~Dong, Y.~Qiao, and C.~C. Loy.
\newblock ESRGAN: Enhanced Super-Resolution Generative Adversarial Networks.
\newblock In {\em ECCV Workshops}, 2018.

\bibitem{heusel2017gans}
M.~Heusel, H.~Ramsauer, T.~Unterthiner, B.~Nessler, and S.~Hochreiter.
\newblock GANs Trained by a Two Time-Scale Update Rule Converge to a Local Nash
Equilibrium.
\newblock In {\em NeurIPS}, 2017.

\bibitem{zhang2018unreasonable}
R.~Zhang, P.~Isola, A.~A. Efros, E.~Shechtman, and O.~Wang.
\newblock The Unreasonable Effectiveness of Deep Features as a Perceptual Metric.
\newblock In {\em CVPR}, 2018.
\end{thebibliography}

\end{document}
